\begin{exercise}{120}
Consider the simplicial set $S=\Delta^2/\Delta^{\{0,1\}}$ and the canonical map
$p\colon\allowbreak\Delta^{\{0,2\}}\to S$. Show that the diagram
\[
\begin{tikzcd}
 \Lambda^2_0 \arrow[r,"f"] \arrow[d] & \Delta^{\{0,2\}} \arrow[d] \\
 \Delta^2 \arrow[r,"c"'] & S
\end{tikzcd}
\]
is a pullback, where $c$ is the canonical projection. Then proceed to show:
\begin{enumerate}[label=(\arabic*)]
 \item $p$ is an inner fibration.
 \item The edge $f|_{\Delta^{\{0,1\}}}$ is degenerate, but not $p$-cocartesian.
\end{enumerate}

\begin{proof}
Compute the quotient degreewise. It collapses all nondecreasing words
on $\{0,1\}$ to one simplex $*$ and leaves every word containing $2$ distinct:
\[
 S_n=\{*\}\amalg\{\alpha\colon[n]\to[2]\text{ monotone}\mid
                         2\in\operatorname{im}\alpha\}.
\]
In particular, $p$ is a monomorphism.
A word in $(\Delta^2)_n$ maps into its image precisely when it avoids
$1$ or avoids $2$.
Thus $c^{-1}(\operatorname{im}p)=\Delta^{\{0,1\}}\cup\Delta^{\{0,2\}}
=\Lambda^2_0$, proving the pullback assertion.
The resulting map $f$ is constant at $0$ on $\Delta^{\{0,1\}}$
and the identity on $\Delta^{\{0,2\}}$.

(1) Consider an inner-horn lifting problem for $p$, with bottom simplex
$\tau\colon\Delta^n\to S$.
If $\tau=*$, lift it constantly to $0$.
Otherwise $\tau$ is represented by a unique nondecreasing word $\alpha$
containing $2$.
If it also contains $1$, some consecutive values are $1,2$.
Every inner horn contains the spine, so the prescribed horn would map
one edge to $c(1,2)$, which is outside the image of $p$; this is impossible.
Hence $\alpha$ takes values in $\{0,2\}$ and itself gives a lift.
In both cases the lift agrees with the prescribed horn because $p$ is a monomorphism.

(2) The description of $f$ makes the indicated edge degenerate.
The displayed square is an outer-horn lifting problem with this marked first edge.
It has no diagonal: such a lift would put $c(0,1,2)$ in the image of $p_2$,
contrary to the degreewise description above.
Definition~3.1.1 therefore says that this edge is not $p$-cocartesian.
\end{proof}
\end{exercise}
